% Chapter 4 outline. Populate only after the methods in Chapter 3 have been
% executed. Report unfavorable, null, and unexpected results alongside expected
% findings; do not write anticipated results as completed evidence.

\newpage
\section{Results and Discussion}
\label{Results and Discussion}
% Organize results in the same order as the approved research questions and
% objectives. For each result: present evidence, interpret it, compare it with
% related literature/systems, and explain its practical meaning and limitations.

\subsection{System Development and Scenario-Analysis Results}
% Report evidence on:
% - final study area, partner office, users, and deployment environment;
% - acquired datasets, provenance, completeness, cleaning, and limitations;
% - implemented architecture, modules, and data flow;
% - implemented routing, center-allocation, time-estimation, resource, and flood
%   or hazard-scenario methods;
% - representative scenario inputs and generated map/decision outputs;
% - functional, integration, performance, and scenario-test outcomes;
% - how blocked roads and full centers change recommendations;
% - traceability from requirements to completed capabilities.
% Add verified maps, diagrams, tables, and screenshots with descriptive captions.

\subsection{System Evaluation Results}
% Report evidence using the approved evaluation plan:
% - participant/validator characteristics without exposing private information;
% - instrument reliability and data-quality checks;
% - route, allocation, time, performance, usability, or decision-support metrics;
% - comparison with the approved baseline or expert reference;
% - findings for each research question and objective;
% - practical implications for the partner organization;
% - threats to validity, dataset limitations, and generalizability.
% Do not claim effectiveness beyond the measured study conditions.
